\documentclass[10pt,letterpaper,oneside]{article}
\usepackage[T1]{fontenc}
\usepackage[utf8]{inputenc}
\usepackage{lmodern}
\usepackage{geometry}
\geometry{margin=0.84in,headheight=21pt,footskip=28pt}
\usepackage{microtype}
\usepackage{xcolor}
\definecolor{ink}{HTML}{15253D}
\definecolor{accent}{HTML}{086A84}
\definecolor{pale}{HTML}{EEF5F8}
\definecolor{muted}{HTML}{4E6474}
\definecolor{grayline}{HTML}{D7E2E8}
\usepackage{graphicx}
\usepackage{fancyhdr}
\usepackage{titlesec}
\usepackage{enumitem}
\usepackage{array,tabularx,longtable,booktabs}
\usepackage[most]{tcolorbox}
\usepackage{xurl}
\usepackage{hyperref}
\usepackage{listings}
\usepackage{amsmath,amssymb}
\usepackage{tikz}
\usetikzlibrary{arrows.meta,positioning,shapes.geometric,calc}
\hypersetup{colorlinks=true,linkcolor=accent,urlcolor=accent,pdftitle={Cryptography, Identity and Privacy Engineering},pdfauthor={Independent technical engineering study companion}}
\renewcommand{\familydefault}{\sfdefault}
\linespread{1.07}
\setlength{\parindent}{0pt}
\setlength{\parskip}{5.5pt}
\emergencystretch=2em
\setlist[itemize]{leftmargin=1.5em,itemsep=2.5pt,topsep=3pt}
\setlist[enumerate]{leftmargin=1.7em,itemsep=3pt,topsep=3pt}
\titleformat{\section}{\Large\bfseries\color{ink}}{\thesection}{0.65em}{}[\vspace{-2pt}{\color{accent}\titlerule}]
\titleformat{\subsection}{\large\bfseries\color{accent}}{\thesubsection}{0.6em}{}
\titleformat{\subsubsection}{\bfseries\color{ink}}{\thesubsubsection}{0.6em}{}
\titlespacing*{\section}{0pt}{16pt}{9pt}
\titlespacing*{\subsection}{0pt}{11pt}{6pt}
\pagestyle{fancy}\fancyhf{}\fancyhead[L]{\footnotesize\color{muted} COMPUTER \& INFORMATION SECURITY ENGINEERING}\fancyhead[R]{\footnotesize\color{accent} SUITE / 5}\fancyfoot[L]{\footnotesize\color{muted} Independent study companion \quad|\quad 2026}\fancyfoot[R]{\footnotesize\color{muted}\thepage}\renewcommand{\headrulewidth}{0.4pt}\renewcommand{\headrule}{\hbox to\headwidth{\color{grayline}\leaders\hrule height \headrulewidth\hfill}}
\newtcolorbox{sourcebox}{enhanced,breakable,boxrule=0pt,leftrule=3pt,colframe=accent,colback=pale,arc=0pt,left=9pt,right=7pt,top=7pt,bottom=7pt}
\newtcolorbox{warningbox}{enhanced,breakable,boxrule=0.5pt,colframe=grayline,colback=white,arc=1pt,left=9pt,right=9pt,top=7pt,bottom=7pt}
\newcommand{\handbookref}[1]{\textsf{\color{accent}\textbf{[H:#1]}}}
\newcommand{\modernext}{\textsf{\color{muted}\textbf{[E]}}}
\newcommand{\designed}{\textsf{\color{accent}\textbf{[D]}}}
\newcommand{\deliverable}[1]{\textbf{Deliverable:} #1}
\lstdefinestyle{lab}{basicstyle=\ttfamily\scriptsize,keywordstyle=\color{accent}\bfseries,commentstyle=\color{muted},numbers=left,numberstyle=\tiny\color{muted},numbersep=7pt,breaklines=true,showstringspaces=false,frame=single,rulecolor=\color{grayline},backgroundcolor=\color{pale},xleftmargin=1em,framexleftmargin=0.6em,tabsize=2,columns=fullflexible}
\lstset{style=lab}
\renewcommand{\arraystretch}{1.21}
\newcommand{\makereportcover}[4]{%
\thispagestyle{empty}\vspace*{1.15in}
{\small\bfseries\color{accent} COMPUTER AND INFORMATION SECURITY ENGINEERING}\par\vspace{13pt}
{\Huge\bfseries\color{ink} #1\par}\vspace{15pt}
{\Large\color{muted} #2\par}\vspace{18pt}
\color{accent}\rule{0.94\textwidth}{2.1pt}\par\vspace{20pt}
{\normalsize\bfseries Document #3 of 8 \quad\textbullet\quad 16 October 2026\par}
\vspace{0.35in}
\begin{sourcebox}\textbf{Orientation.} #4\end{sourcebox}
\vfill
{\small\color{muted} Independent, original engineering study and practice materials grounded in\par
\emph{Computer and Information Security Handbook}, 4th ed., edited by John R. Vacca,\par
Elsevier/Morgan Kaufmann, 2025, Volumes I--II. Not a reproduction or official standard.}
\newpage\setcounter{page}{1}\tableofcontents\newpage
}
\begin{document}
\makereportcover{Cryptography, Identity and Privacy Engineering}{Data protection, PKI, key lifecycles, MFA, privacy and post-quantum transition}{5}{Connect cryptographic fundamentals and public-key infrastructure to practical system decisions, identity-assurance engineering and privacy controls.}

\section{Threat-model cryptography before choosing algorithms}
The handbook provides a cryptographic primer, block/public-key ciphers, PKI, password-authenticated key establishment and context-aware multifactor authentication. Volume II adds online privacy, privacy-enhancing technology, virtual private networks, advanced encryption and data loss protection. \handbookref{3,4,47--51,53--61,82,97,100}
\begin{sourcebox}\textbf{Decision sequence:} state asset and adversary; choose required confidentiality/integrity/authenticity properties; select a maintained protocol/library; define identity and keys; design lifecycle and recovery; test misuse and failure paths. Cryptography does not repair improper authorization or compromised endpoints.\end{sourcebox}

\section{Primitives and correct composition}
\par\medskip\noindent\begin{tabularx}{\textwidth}{@{}p{0.24\textwidth}X X@{}}\toprule
Primitive & Intended security role & Common misuse \\ \midrule
CSPRNG & Unpredictable keys, nonces, tokens & Predictable time-seeded tokens\\
Hash & Fixed digest for integrity structures & Treating unkeyed hash as authentication\\
HMAC & Message authentication with shared key & Weak secret and ambiguous encoding\\
KDF / password hash & Turn weak human secret into resistant verifier & Fast general-purpose hash for passwords\\
AEAD & Encrypt and authenticate data together & Nonce reuse; missing context binding\\
Signature & Verify origin/integrity under public key & Trusting an unvalidated public key\\
KEM / key agreement & Establish a shared secret & No authentic binding to peer identity\\
\bottomrule\end{tabularx}\par\medskip
\handbookref{3,47,50,82,100} \modernext{} Prefer well-reviewed high-level cryptographic APIs over hand-implemented ciphers or custom protocol combinations. DES and Triple DES discussions in the handbook are historical/educational, not blanket deployment endorsements.
\subsection{Authenticated encryption data model}\modernext{}
With authenticated encryption with associated data (AEAD), an application sends a nonce \(N\), authenticated context \(A\), and encrypted payload \(C\). Conceptually,
\[ (C,T)\gets\operatorname{Encrypt}(K,N,P,A),\qquad P\gets\operatorname{Decrypt}(K,N,C,T,A). \]
The tag \(T\) rejects modified ciphertext or modified associated context. A nonce must follow the \emph{specific algorithm's} uniqueness requirements; secure key and nonce management are application responsibilities. Do not interpret the equation as an implementation specification.

\section{Key and secret lifecycle}
\par\medskip\noindent\begin{longtable}{@{}p{0.18\textwidth}p{0.37\textwidth}p{0.35\textwidth}@{}}\toprule
Stage & Engineering question & Required evidence \\ \midrule\endhead
Generate & Was entropy strong and algorithm approved for purpose? & CSPRNG/KMS creation record\\
Provision & Which workload or person may access key? & Explicit access policy and identity proof\\
Use & Are cryptographic operations and logs controlled? & Invocation logs and configuration test\\
Rotate & Can dependencies switch without data loss? & Dual-read migration and rollback test\\
Revoke & What happens if compromised? & Revocation event and propagation check\\
Recover & Is business recovery feasible without shared-fate failure? & Restore protocol and access separation\\
Retire & Is data retention/cryptographic erasure understood? & Deletion and key-retirement record\\
\bottomrule\end{longtable}\par\medskip
\designed{} Inventory data-at-rest keys, signing keys, TLS credentials, service tokens, SSH keys, device identities and backup keys independently. Classify rotation as scheduled, compromise-triggered or lifecycle-event-triggered; indiscriminate rotation can create outages.

\section{Public key infrastructure and TLS}
\handbookref{49} Public keys are only useful as identity assertions if validated under an appropriate trust model. An X.509 certificate binds a subject name to a key under a signature chain; validity dates, intended usage, revocation information and trust anchor must all be evaluated by the intended client.
\begin{enumerate}
\item Inspect server-provided identity name(s), certificate path and algorithm.
\item Confirm client trust store and application-specific constraints, including hostname validation.
\item Verify expiration handling, revocation strategy and replacement automation.
\item Test rotation without clients silently bypassing validation.
\item Protect private keys and constrain certificate issuance and signing authorities.
\end{enumerate}
\modernext{} TLS 1.3, automated certificate issuance and workload identity are contemporary deployment topics. A secure channel does not prove that the remote service is authorized to disclose a specific object.

\section{Digital identity architecture}
\subsection{Separate proofing, authentication, federation and authorization}
\handbookref{4,51,73,79} An application must not conflate an authenticated identity with an authorized operation. NIST SP 800-63-4 addresses identity proofing, authentication and federation; authorization is a separate application-policy decision. \modernext{} [N4]
\par\medskip\noindent\begin{tabularx}{\textwidth}{@{}p{0.23\textwidth}X X@{}}\toprule
Function & Primary question & Example test \\ \midrule
Identity proofing & Is the claimed identity established to sufficient assurance? & Verify enrollment/identity proofing process\\
Authentication & Is this claimant legitimately controlling authenticator? & Replay, phishing-resistance and recovery tests\\
Federation & Can receiving service trust assertion and issuer? & Audience, issuer, signature and expiry checks\\
Authorization & May this principal access this action/object? & Negative cross-role and cross-tenant cases\\
Session lifecycle & Does previous trust persist appropriately? & Revocation, idle and high-risk action tests\\
\bottomrule\end{tabularx}\par\medskip
\subsection{Password, MFA and recovery controls}\modernext{}
Use an approved password storage scheme rather than encrypting passwords for lookup. Favor phishing-resistant authenticators for privileged operations where feasible; protect enrollment, account recovery and factor reset just as carefully as the login flow. A second step that simply repeats the same compromised channel is not necessarily independent assurance.

\section{Privacy engineering and data minimization}
\handbookref{53--58,60,95,97} Privacy is not equivalent to encrypting stored personal information. A processing system may be confidential yet over-collect, over-retain, profile, or redistribute personal information beyond its legitimate purpose.
\subsection{Data processing inventory}\designed{}
\par\medskip\noindent\begin{tabularx}{\textwidth}{@{}p{0.18\textwidth}X@{}}\toprule
Field & Example for usage analytics \\ \midrule
Purpose & Diagnose service reliability; avoid unrelated marketing use\\
Categories & Pseudonymous session identifier and event timestamp\\
Collection & Application edge events with sensitive fields excluded\\
Retention & Defined operational window; automatic deletion policy\\
Recipients & Named operations processor under approved transfer terms\\
Protection & Role-scoped access, masking and audit logs\\
Rights/requests & Search, export and deletion workflow where applicable\\
Verification & Payload inspection and retention-expiry test\\
\bottomrule\end{tabularx}\par\medskip
All privacy obligations depend on law, jurisdiction, contracts and processing context; this inventory is a design aid, not legal advice.

\section{Network confidentiality and secure communication}
\handbookref{13,48,52,59,61} VPNs, messaging systems, voice protocols and satellite links involve different endpoint, authentication and traffic-analysis threats. Evaluate endpoints and key lifecycle alongside tunnel or transport encryption. Encryption hides content more readily than all metadata; trust boundaries and retention still matter.

\section{Post-quantum cryptography migration}\modernext{}
NIST finalized FIPS 203 (ML-KEM), FIPS 204 (ML-DSA), and FIPS 205 (SLH-DSA) in 2024. These standards do not mean that every legacy protocol can instantly replace a cipher without ecosystem work. [N5]
\begin{enumerate}
\item Inventory asymmetric cryptography and long-lived secrets by system and business owner.
\item Separate encryption/key-establishment exposure from signing/certificate exposure.
\item Assess ``harvest now, decrypt later'' sensitivity and cryptoperiod.
\item Prioritize supported protocol/library upgrades and vendor dependency requirements.
\item Test interoperability, message sizes, performance, failure handling and recovery.
\item Document transition milestones and exceptions; retain algorithm agility.
\end{enumerate}

\section{Worked code: safe library use}\designed{}
This example requires the separately installed \texttt{cryptography} Python package; it is a learning illustration, not a substitute for production key management.
\begin{lstlisting}[language=Python,caption={AES-GCM with associated tenant context}]
import os
from cryptography.hazmat.primitives.ciphers.aead import AESGCM
key = AESGCM.generate_key(bit_length=256)
aead = AESGCM(key)
nonce = os.urandom(12)  # never reuse with same key
context = b"tenant:example|record:42"
message = b"demo data only"
ciphertext = aead.encrypt(nonce, message, context)
restored = aead.decrypt(nonce, ciphertext, context)
assert restored == message
# Production: managed keys, nonce strategy, rotation, and storage.
\end{lstlisting}
\section{Security review exercises}
Identify five ways valid encryption can fail to protect sensitive data; threat-model factor reset; review an untrusted certificate; produce a key inventory/rotation plan for three services; plan phased post-quantum transition for high-value, long-lived data. Assessment: correct threat boundaries (30), key lifecycle (30), identity and authorization (20), privacy and recovery (20). \designed{}

\section*{Handbook chapter alignment}\addcontentsline{toc}{section}{Handbook chapter alignment}
These chapters ground the domain scope. The document offers original synthesis and technical procedures, not a chapter-by-chapter reproduction.
\par\medskip\noindent\begin{longtable}{@{}p{0.085\textwidth}p{0.68\textwidth}p{0.16\textwidth}@{}}\toprule Chapter & Handbook topic & Combined PDF p. \\ \midrule\endhead3 & A Cryptography Primer & 59 \\
4 & Verifying User and Host Identity & 83 \\
13 & Internet Security & 271 \\
47 & Data Encryption & 801 \\
48 & Satellite Encryption & 827 \\
49 & Public Key Infrastructure & 841 \\
50 & Password-Based Authenticated Key Establishment Protocols & 863 \\
51 & Context-Aware Multifactor Authentication Survey & 879 \\
52 & Instant-Messaging Security & 891 \\
53 & Online Privacy & 957 \\
54 & Privacy-Enhancing Technologies & 977 \\
55 & Personal Privacy Policies & 993 \\
56 & Detection of Conflicts in Security Policies & 1003 \\
57 & Supporting User Privacy Preferences in Digital Interactions & 1023 \\
58 & Privacy and Security in Environmental Monitoring Systems: Issues and Solutions & 1045 \\
59 & Virtual Private Networks & 1065 \\
60 & Identity Theft & 1079 \\
61 & VoIP Security & 1093 \\
72 & Biometrics & 1247 \\
73 & Online Identity and User Management Services & 1265 \\
79 & Access Controls & 1353 \\
82 & Fundamentals of Cryptography & 1377 \\
95 & Online e-Reputation Management Services & 1595 \\
97 & Data Loss Protection & 1629 \\
100 & Advanced Data Encryption & 1691 \\
\bottomrule\end{longtable}\par\medskip

\section*{Source and reference register}
\addcontentsline{toc}{section}{Source and reference register}
\begin{description}[leftmargin=1.6em,style=nextline]
\item[H] John R. Vacca (ed.), \emph{Computer and Information Security Handbook}, 4th ed., Volumes I--II, Elsevier/Morgan Kaufmann, 2025. Chapter identifiers in this document point to this book.
\item[N1] NIST, \emph{Cybersecurity Framework 2.0} (2024), \url{https://www.nist.gov/cyberframework}.
\item[N2] NIST SP 800-218, \emph{Secure Software Development Framework v1.1} (2022), \url{https://csrc.nist.gov/pubs/sp/800/218/final}. Later drafts are not treated as final requirements.
\item[N3] NIST SP 800-61 Rev. 3, \emph{Incident Response Recommendations and Considerations for Cybersecurity Risk Management} (2025), \url{https://csrc.nist.gov/pubs/sp/800/61/r3/final}.
\item[N4] NIST SP 800-63-4, \emph{Digital Identity Guidelines} (2025), \url{https://csrc.nist.gov/pubs/sp/800/63/4/final}.
\item[N5] NIST, FIPS 203/204/205 (2024), \url{https://csrc.nist.gov/Projects/Post-Quantum-Cryptography/publications}.
\item[O1] OWASP, \emph{API Security Top 10} (2023), \url{https://api-security.owasp.org/editions/2023/en/0x11-t10/}.
\item[C1] CISA, \emph{Known Exploited Vulnerabilities Catalog}, \url{https://www.cisa.gov/known-exploited-vulnerabilities-catalog}.
\item[M1] MITRE ATT\&CK Enterprise, \url{https://attack.mitre.org/matrices/enterprise/}.
\item[G1] GitHub, \emph{REST API endpoints for code scanning}, \url{https://docs.github.com/en/rest/code-scanning}.
\end{description}
\begin{warningbox}\textbf{Interpretation rule.} \handbookref{number} identifies a conceptual anchor in the supplied handbook, not a verbatim quotation or an official requirement. \modernext{} identifies a contemporary extension informed by external references or current engineering practice. \designed{} identifies an independently developed example, template, exercise, metric, or operational procedure. Requirements are binding only if separately adopted through applicable law, contract, policy, or the cited standards; this suite does not assert otherwise.\end{warningbox}

\end{document}
